\section{The nine variable ratio interface}
We first use exactly the mixed-domain coordinates
\[
 b,g,h,s,w,i\in\Pos,\qquad \phi,j,n\in\N.
\]
Every other symbol is a polynomial abbreviation. The source registers
$R,N,G_c=G_{\rm cap}$ are those of Section~\ref{sec:carry8}, with the
selected even- or odd-$K_c$ threshold formula. They satisfy the imported
ratio theorem's unconditional size requirements before a zero is assumed.
The following quotient removal and conditional norm lemma are retained
from the preceding iteration; the new source and sharper majorant are
composed with their full two-way proofs, not merely with sample tests.

\subsection{Literal ratio registers and exact quotient removal}
Define
\begin{align*}
 U&=N^2w,&Y&=N^2s,&M_0&=RY,\\
 P&=2M_0^2U,&K&=R+1+h(P-1),&A&=M_0(U+1),\\
 B_1&=2R+1,&W_1&=bw,&C&=B_1+W_1+\phi,\\
 J&=4A-5,&Q&=3W_1C-2(W_1^2-1),&&\\
 D&=(A^2-1)C^2+1,&k&=iJC,&L&=kD,\\
 E&=LC,&F&=(A^2-1)E^2+1,&&\\
 G&=A+F(CD+1-A),&H&=C+F(B_1+jC),&&\\
 I&=(G^2-1)H^2+1,&a_0&=(P^2-1)K^2+1,&a_1&=DFI,\\
 m&=[K^2-4(C-KY)^2]G_c-8gK^2.&&&&
\end{align*}
The inherited ratio theorem instead has $H_{\rm old}=B_1+j_{\rm old}C$,
$j_{\rm old}\ge1$, and divisibility
$JF\mid QF+J(H_{\rm old}-C)$. Since $F\equiv1\pmod J$ and
$F\equiv1\pmod C$, $\gcd(F,J)=1$, so that divisibility is equivalent to
\[
 J\mid Q\quad\hbox{and}\quad F\mid H_{\rm old}-C.
\]
If the second holds, $z=(H_{\rm old}-C)/F$ is positive and congruent to
$B_1$ modulo $C$. As $0<B_1<C$, uniquely $z=B_1+jC$ for $j\ge0$.
This is exactly the new $H$. Conversely every such new $j$ gives the old
positive integer
\[
 j_{\rm old}=1+\frac{F-1}{C}B_1+Fj.
\]
The quotient is an integer because $C\mid F-1$. This is a bijection of the
old divisible witnesses with the new unrestricted nonnegative witnesses,
preserving $H$ itself. No cofinal-witness assertion or extra hypothesis is
introduced. The resulting ratio conditions are precisely
\[
 a_0\text{ square},\quad a_1\text{ square},\quad J\mid Q,\quad m>0.
\]
The numerical Jones theorem gives from these conditions that $b$ is a power
of two, $8g<G_c$, and $N^2\mid\binom{2R}{R}$. Conversely the stronger
source bound $32g<G_c$ produces the requisite ratio witnesses.

\subsection{A general conditional-separation norm theorem}
\begin{theorem}\label{thm:conditional}
Let $u,v,m,a_0,a_1,V_0,V_1$ be integers with $u>0$,
$0<v/u<V_0+V_1$, $a_i>0$, $V_i>0$, $a_i<V_i^2$, and
$m>0\Rightarrow a_0\ne a_1$. Put
\[
 \Omega=um,\qquad \Theta=un+v-\Omega(V_0+V_1).
\]
Then there exists $n\in\N$ satisfying
\[
 (\Theta^2-\Omega^2(a_0+a_1))^2-4\Omega^4a_0a_1=0
\]
if and only if $m>0$, $a_0,a_1$ are integer squares, and $u\mid v$.
\end{theorem}
\begin{proof}
If $m=0$, $\Theta=un+v>0$, excluding a zero. Otherwise factor the norm
over the reals to obtain signed roots $r_i^2=a_i$ with
$\Theta/\Omega=r_0+r_1$. Rearrangement gives
\[
 n+v/u=mZ,\qquad Z=V_0+V_1+r_0+r_1>0.
\]
The left side is positive, so $m>0$ \emph{before} integrality of either
root is known. Conditional separation now makes $x=r_0+r_1$ nonzero;
otherwise $a_0=a_1$. Since $x$ is rational,
$2xr_0=x^2+a_0-a_1$ shows $r_0,r_1\in\mathbb Q$.
A rational number whose square is an integer is an integer, as follows
by writing it as a reduced fraction. Hence $Z\in\Z$ and $u\mid v$.
Conversely choose the positive integer roots and put $n=mZ-v/u$.
Then $n\in\Z$ and $n\ge Z-v/u>V_0+V_1-v/u>0$.
Substitution makes a linear factor zero. The quotient bound is needed only
for completeness, not for the positivity-first soundness argument.
\end{proof}
Positive majorants are essential, not just $V_i^2>a_i$. Without separation,
$a_0=a_1=2$, $V_0=V_1=2$, $u=v=2$, $m=1$, $n=3$ is a false square
witness. Without positive majorants, $a_0=1,a_1=4,V_0=-2,V_1=4$,
$u=v=1,m=-1,n=0$ is a false positive-margin witness, even though the
squared-majorant and quotient inequalities hold. These are retained
as negative controls.

\subsection{The lower cost simple majorant companion}
Unconditionally $A\ge2$ and $P,K,C,J,D,E,F,G,H>1$. Define
\[
 V_0=PK,\qquad T=k(8D^2-4D-1),\qquad V_1=TGH.
\]
The identities $E=CL$, $L=kD$ give $F=(D-1)k^2D^2+1$. Direct expansion
then yields
\[
 V_0^2-a_0=K^2-1>0,\qquad
 T^2-64DF=(8D+1)k^2-64D=8D(k^2-8)+k^2>0,
\]
because $k=iJC\ge3$. Therefore the strictly positive majorant $V_1$ has
\[
 \Delta_1=V_1^2-64a_1
 =[(8D+1)k^2-64D](GH)^2+64DF(H^2-1)>0.
\]
The lower degree of this polynomial defect, rather than simply a smaller
numerical bound, is the degree improvement.

If $m>0$, its defining margin implies $4(C-KY)^2<K^2$, so
$C>K(Y-1/2)>K$ because $Y\ge64$. Also
\[
 A^2-P=M_0^2(U^2+1)>0,\quad D>C^2,\quad F\ge A^2,\quad
 G\ge A,\quad I\ge G^2\ge A^2.
\]
For $G\ge A$, use $G=A+F(CD+1-A)$ and $D\ge A^2$, $C\ge1$.
It follows that
\[
 a_1=DFI>C^2A^4>C^2P^2>K^2P^2>a_0.
\]
This separation is established without assuming squarehood or divisibility.
The same positivity-first proof works without the auxiliary quotient offsets.
Indeed $Q<3C^2$, $D>C^2$ and $T>3D$ imply $V_1>Q$.
The lower-cost simple-majorant fixed polynomial is
\begin{align*}
 \Omega&=Jm,\\
 \Theta&=Jn+Q-\Omega(V_0+V_1),\\
 F_9&=[\Theta^2-\Omega^2(a_0+64a_1)]^2
                   -256\Omega^4a_0a_1=0.
\end{align*}
To check the offset-free form explicitly, $m=0$ gives $\Theta=Jn+Q>0$.
Otherwise a zero factor gives $n+Q/J=mZ$ with
$Z=V_0+V_1+r_0+r_1>0$. Thus $m>0$, and the separation just proved
makes $r_0+r_1$ nonzero. The rational reconstruction from
Theorem~\ref{thm:conditional} gives integer roots and $J\mid Q$.
Conversely choose the positive integer roots and set $n=mZ-Q/J$.
As $Z>V_1>Q>Q/J$, this is a positive integer and a factor vanishes.

Squarehood of $64a_1$ is equivalent to squarehood of $a_1$: an integer
square root divided by 8 is a rational square root of an integer, hence
an integer. The finalizer is therefore equivalent to exactly the four
ratio conditions. Composing both directions proves the desired source
solvability equivalence.

For the positive-domain polynomial, replace $\phi,j,n$ in this mixed-domain
formula by the corresponding positive coordinates minus 1. For the
nonnegative-domain polynomial, replace $b,g,h,s,w,i$ by the corresponding
nonnegative coordinates plus 1. These are bijections; the supplied circuit
ledgers charge their operations. They must not be omitted when stating the
domain.

